\documentclass[11pt,a4paper]{article}

\def\nyear {2026}
\def\nterm {Winter}
\def\nlecturer {Howard Nuer and notes by Lev Borisov}
\def\ncourse {Algebraic Geometry}

\makeatletter

% packages
\usepackage{amssymb,amsfonts,amsmath,calc,tikz,pgfplots,geometry,mathtools}
\usepackage{color}   % May be necessary if you want to color links
\usepackage[svgnames]{xcolor}
\usepackage[hidelinks]{hyperref}
\usepackage{forest}
\usepackage{commath} % ew
\usepackage{esvect} % \vv makes vectors
\usepackage{centernot} % for the comparison
\usepackage{esdiff}
\usepackage{esint}
\usepackage{amsthm}
\usepackage{fancyhdr}
\usepackage{bm}
\usepackage{witharrows}
\usepackage{bookmark}
\usepackage{tikz-cd}
\usepackage{quiver}
\usepackage{bbm}
\usepackage{textcomp}
\usepackage{float}
\usepackage{gensymb}
\usepackage{cleveref}
\usepackage{environ}
\usepackage{comment}
\usepackage{cancel}
\usepackage{xparse}

% Inevitable cs
\usepackage{listings}
\usepackage[]{algorithm2e}

% tikz libraries
\usetikzlibrary{positioning}
\usetikzlibrary{matrix}
\usetikzlibrary{arrows}
\usetikzlibrary{bending}
\usetikzlibrary{arrows.meta}
\usetikzlibrary{decorations.markings}
\usetikzlibrary{decorations.pathreplacing}
\usetikzlibrary{decorations.markings}
\usetikzlibrary{patterns}
\usetikzlibrary{backgrounds}

% Page style setup
\pagestyle{fancy}
\fancyhfoffset{0pt}
\renewcommand{\headwidth}{\textwidth}
\geometry{margin=1in}
\pgfplotsset{compat=1.18}
\setlength{\headheight}{14.6pt}
\addtolength{\topmargin}{-1.6pt}
\hypersetup{
    colorlinks=false,
    linktoc=section,
    linkcolor=black,
}

%% maketitle setup
\ifx \nauthor\undefined
  \def\nauthor{Yeheli Fomberg}
\else
\fi

\ifx \nemail\undefined
  \def\nemail{\href{mailto:yp@campus.technion.ac.il}
  {\texttt{<yp@campus.technion.ac.il>}}}
\else
\fi

\ifx \ncoursehead \undefined
  \def\ncoursehead{\ncourse}
\fi

\lhead{\emph{\nouppercase{\leftmark}}}
\ifx \nextra \undefined
  \rhead{
    \ifnum\thepage=1
    \else
      \ncoursehead
    \fi}
\else
  \rhead{
    \ifnum\thepage=1
    \else
      \ncoursehead \ (\nextra)
    \fi}
\fi

\def\notes{notes}
\def\problems{problems}

\ifx \ntype \undefined
  \def\ntype{notes}
\else
\fi

\ifx \ntype \notes
  \let\@real@maketitle\maketitle
  \renewcommand{\maketitle}{\@real@maketitle\begin{center}
  \begin{minipage}[c]{0.9\textwidth}\centering\footnotesize
  These notes are not endorsed by the lecturers.
  I have revised them outside lectures to incorporate supplementary
  explanations, clarifications, and material for fun.
  While I have strived for accuracy, any errors or misinterpretations 
  are most likely mine.
  \end{minipage}\end{center}}
\else
  \ifx \ntype \problems
    \let\@real@maketitle\maketitle
    \renewcommand{\maketitle}{\@real@maketitle\begin{center}
    \begin{minipage}[c]{0.9\textwidth}\centering\footnotesize
    These problems are mostly based on the problem sets I was given
    when taking the course.
    
    While I have strived for accuracy and completeness,
    some of the problems would not have been here without the help of
    friends.
    \end{minipage}\end{center}}
  \else
  \fi
\fi


% theorem environments
\theoremstyle{definition}
\newtheorem{definition}{Definition}[section]
\newtheorem{remark}{Remark}[section]
\newtheorem{example}{Example}[section]
\newtheorem{exercise}{Exercise}[section]
\newtheorem{paradox}{Paradox}[section]
\newtheorem{entry}{Entry}[section]
\newtheorem*{solution}{Solution}
\newtheorem*{question}{Question}
\newtheorem*{answer}{Answer}
\newtheorem*{problem}{Problem}
\theoremstyle{plain}
\newtheorem{theorem}{Theorem}[section]
\newtheorem{proposition}[theorem]{Proposition}
\newtheorem{lemma}[theorem]{Lemma}
\newtheorem{corollary}[theorem]{Corollary}
\newenvironment{proofof}[1]{%
    \renewcommand{\proofname}{Proof of #1}%
    \begin{proof}
}{%
    \end{proof}
}
% tikz customization
\pgfarrowsdeclarecombine{twolatex'}{twolatex'}{latex'}{latex'}{latex'}{latex'}
\tikzset{->/.style = {decoration={markings,
                                  mark=at position 0.999
                                  with {\arrow[scale=2]{latex'}}},
                      postaction={decorate}}}
\tikzset{<-/.style = {decoration={markings,
                                  mark=at position 0.001 with {\arrowreversed[scale=2]{latex'}}},
                      postaction={decorate}}}
\tikzset{-->/.style = {decoration={markings,
                                  mark=at position 0.999
                                  with {\arrow[scale=2]{latex'[sep=0pt -2.5]}}},
                      postaction={decorate}}}
\tikzset{<--/.style = {decoration={markings,
                                  mark=at position 0.001
                                  with {\arrowreversed[scale=2]{latex'[sep=0pt -2.5]}}},
                      postaction={decorate}}}
\tikzset{<->/.style = {decoration={markings,
                                  mark=at position 0.001 with {\arrowreversed[scale=2]{latex'}},
                                  mark=at position 0.999 with {\arrow[scale=2]{latex'}},},
                      postaction={decorate}}}
\tikzset{->-/.style = {decoration={markings,
    mark=at position 0.50+0.225/#1 with {\arrow[scale=2]{latex'}}},
                      postaction={decorate}}}
\tikzset{-<-/.style = {decoration={markings,
                                   mark=at position 0.50-0.225/#1 with {\arrowreversed[scale=2]{latex'}}},
                       postaction={decorate}}}
\tikzset{->>/.style = {decoration={markings,
                                  mark=at position 1 with {\arrow[scale=2]{latex'}}},
                      postaction={decorate}}}
\tikzset{<<-/.style = {decoration={markings,
                                  mark=at position 0 with {\arrowreversed[scale=2]{twolatex'}}},
                      postaction={decorate}}}
\tikzset{<<->>/.style = {decoration={markings,
                                   mark=at position 0 with {\arrowreversed[scale=2]{twolatex'}},
                                   mark=at position 1 with {\arrow[scale=2]{twolatex'}}},
                       postaction={decorate}}}
\tikzset{->>-/.style = {decoration={markings,
                                   mark=at position 0.50+0.450/#1 with {\arrow[scale=2]{twolatex'}}},
                       postaction={decorate}}}
\tikzset{-<<-/.style = {decoration={markings,
                                   mark=at position 0.50-0.450/#1 with {\arrowreversed[scale=2]{twolatex'}}},
                       postaction={decorate}}}

\pgfarrowsdeclare{biggertip}{biggertip}{%
  \setlength{\arrowsize}{1pt}
  \addtolength{\arrowsize}{.1\pgflinewidth}
  \pgfarrowsrightextend{0}
  \pgfarrowsleftextend{-5\arrowsize}
}{%
  \setlength{\arrowsize}{1pt}
  \addtolength{\arrowsize}{.1\pgflinewidth}
  \pgfpathmoveto{\pgfpoint{-5\arrowsize}{4\arrowsize}}
  \pgfpathlineto{\pgfpointorigin}
  \pgfpathlineto{\pgfpoint{-5\arrowsize}{-4\arrowsize}}
  \pgfusepathqstroke
}
\tikzset{
	EdgeStyle/.style = {>=biggertip}
}

\tikzset{circ/.style = {fill, circle, inner sep = 0, minimum size = 3}}
\tikzset{scirc/.style = {fill, circle, inner sep = 0, minimum size = 1.5}}
\tikzset{mstate/.style={circle, draw, black, text=black, minimum width=0.7cm}}

\tikzset{eqpic/.style={baseline={([yshift=-.5ex]current bounding box.center)}}}

\definecolor{mblue}{rgb}{0.2, 0.3, 0.8}
\definecolor{morange}{rgb}{1, 0.5, 0}
\definecolor{mgreen}{rgb}{0, 0.4, 0.2}
\definecolor{mred}{rgb}{0.5, 0, 0}

% topology
\DeclareMathOperator{\dfr}{dfr} % deformation retract
\newcommand{\Cells}{\text{Cells}}
\newcommand{\RP}{\mathbb{RP}}
\newcommand{\CP}{\mathbb{CP}}

% order theory
\newcommand{\join}{\vee}
\newcommand{\meet}{\wedge}

% algebraic geometry
\DeclareMathOperator{\Rad}{Rad} % Radical of modules
\DeclareMathOperator{\Spec}{Spec}
% \renewcommand{\div}{\mathrm{div}} % divisors
\DeclareMathOperator{\Mer}{Mer} % Meromorphic functions
\DeclareMathOperator{\Exc}{Exc} % Exceptional divisor
\DeclareMathOperator{\Td}{Td} % Todd classes

% algebra

\DeclareMathOperator{\trdeg}{tr.deg.}
\DeclareMathOperator{\odd}{odd}
\DeclareMathOperator{\even}{even}
\DeclareMathOperator{\lcm}{lcm}
\DeclareMathOperator{\ord}{ord}
\DeclareMathOperator{\codim}{codim}
\DeclareMathOperator{\Ann}{Ann}
\DeclareMathOperator{\Ext}{Ext}
\DeclareMathOperator{\Sym}{Sym}
\DeclareMathOperator{\Out}{Out}
\DeclareMathOperator{\Aut}{Aut}
\DeclareMathOperator{\End}{End}
\DeclareMathOperator{\Inn}{Inn}
\DeclareMathOperator{\Mat}{Mat}
\DeclareMathOperator{\Fix}{Fix}
\DeclareMathOperator{\std}{std}
\DeclareMathOperator{\sgn}{sgn}
\DeclareMathOperator{\adj}{adj}
\DeclareMathOperator{\id}{id}
\DeclareMathOperator{\op}{op}
\DeclareMathOperator{\tr}{tr}
\DeclareMathOperator{\diag}{diag}
\DeclareMathOperator{\rank}{rank}
\DeclareMathOperator{\rk}{rk}
\DeclareMathOperator{\coker}{coker}
\DeclareMathOperator{\Mult}{Mult} % multiplier algebra
\DeclareMathOperator{\Frac}{Frac} % fraction field
\DeclareMathOperator{\Rat}{Rat}
\DeclareMathOperator{\Gal}{Gal}
\newcommand{\ioplus}{\overset{\text{int}}{\oplus}} % interior direct sum
\newcommand{\GG}{\mathcal G} % Galois correspondence
\newcommand{\FF}{\mathcal F} % Galois correspondence
\newcommand{\cupp}{\smile} % cup product
\newcommand{\capp}{\frown} % cap product
\newcommand{\actson}{\curvearrowright}
\newcommand{\bra}{\left\langle}
\newcommand{\ket}{\right\rangle}
\newcommand{\idealin}{\triangleleft\,}
\newcommand{\normalin}{\triangleleft\,}
\newcommand{\ip}[2]{\left\langle #1, #2 \right\rangle}
\newcommand{\bigslant}[2]
{{\raisebox{.2em}{$#1$}\left/\raisebox{-.2em}{$#2$}\right.}}
\newcommand{\properideal}{%
  \mathrel{\ooalign{$\lneq$\cr\raise.22ex\hbox{$\lhd$}\cr}}}

% categories
\newcommand{\cat}[1]{\mathbf{#1}}
\newcommand{\Ch}{\cat{Ch}} % Chain complexes   
\newcommand{\Set}{\cat{Set}}
\newcommand{\Grp}{\cat{Grp}}
\newcommand{\Ab}{\cat{Ab}}
\newcommand{\Ring}{\cat{Ring}}
\newcommand{\CRing}{\cat{CRing}}
\newcommand{\Top}{\cat{Top}}
\newcommand{\Vect}{\cat{Vect}}
\newcommand{\cmod}{\cat{mod}}

%% matrix groups
\newcommand{\GL}{\mathrm{GL}}
\newcommand{\Or}{\mathrm{O}}
\newcommand{\PGL}{\mathrm{PGL}}
\newcommand{\PSL}{\mathrm{PSL}}
\newcommand{\PSO}{\mathrm{PSO}}
\newcommand{\PSU}{\mathrm{PSU}}
\newcommand{\SL}{\mathrm{SL}}
\newcommand{\msl}{\mathrm{sl}}
\newcommand{\SO}{\mathrm{SO}}
\newcommand{\Spin}{\mathrm{Spin}}
\newcommand{\Sp}{\mathrm{Sp}}
\newcommand{\SU}{\mathrm{SU}}
\newcommand{\U}{\mathrm{U}}
\let\char\relax
\newcommand{\char}{\text{char}}

% analysis
\renewcommand{\Re}{\mathrm{Re}}
\renewcommand{\Im}{\mathrm{Im}}
\NewDocumentCommand{\dx}{o}{
  \IfNoValueTF{#1}
    {\dif x}                  % if no optional argument
    {\dif^{#1} \!x}         % if optional argument
}
\NewDocumentCommand{\dy}{o}{
  \IfNoValueTF{#1}
    {\dif y}                  % if no optional argument
    {\dif^{#1} \!y}         % if optional argument
}
\newcommand{\dt}{\dif t}
\newcommand{\du}{\dif u}
\newcommand{\dv}{\dif v}
\newcommand{\dz}{\dif z}
\newcommand{\ds}{\dif s}
\newcommand{\dw}{\dif w}
\newcommand{\dr}{\dif r}
\newcommand{\dm}{\dif m}
\newcommand{\df}{\dif f}
\newcommand{\dV}{\dif V}
\newcommand{\dS}{\dif S}
\newcommand{\dA}{\dif \mathrm{A}}
\newcommand{\dmu}{\dif \mu}
\newcommand{\dnu}{\dif \nu}
\newcommand{\dtheta}{\dif \theta}
\newcommand{\domega}{\dif \omega}
\newcommand{\dsigma}{\dif \sigma}
\newcommand{\dtau}{\dif \tau}
\newcommand{\dvarphi}{\dif \varphi}
\renewcommand{\dif}{\mathrm{d}}
\renewcommand{\Dif}{\mathrm{D}}
\renewcommand{\triangle}{\Delta}
\newcommand{\ptriangle}{\partial \triangle}
\DeclareMathOperator{\im}{im}
\DeclareMathOperator{\cis}{cis}
\DeclareMathOperator{\rot}{rot}
\DeclareMathOperator{\vspan}{span}
\DeclareMathOperator{\grad}{grad}
\DeclareMathOperator{\curl}{curl}
\DeclareMathOperator{\esssup}{esssup}
\newcommand{\hodge}{{\mathnormal{\star}}}
\DeclareMathOperator{\range}{range}
\DeclareMathOperator{\Hom}{Hom}
\DeclareMathOperator{\Int}{Int}
\DeclareMathOperator{\Conv}{Conv} % convex hull
\newcommand{\ind}{\mathbbm{1}}
\DeclareMathOperator{\Res}{Res}
\DeclareMathOperator{\graph}{graph}
\DeclareMathOperator{\diam}{diam}
\DeclareMathOperator{\supp}{supp}
\DeclareMathOperator{\Vol}{Vol} % Volume
\DeclareMathOperator{\Area}{Area} % Area
\DeclareMathOperator{\area}{area}
\DeclareMathOperator{\Ind}{Ind} % Index
\newcommand{\length}{\mathrm{length}}

% logic
\DeclareMathOperator{\MOD}{MOD}
\DeclareMathOperator{\Theory}{Theory}
\newcommand{\notimplies}{%
  \mathrel{{\ooalign{\hidewidth$\not\phantom{=}$\hidewidth\cr$\implies$}}}}

% nice
\newcommand{\half}{\frac{1}{2}}
\newcommand{\pair}{\del}
\newcommand{\taking}[1]{\xrightarrow{#1}}
\newcommand{\otaking}[1]{\xleftarrow{#1}}
\newcommand{\inv}{^{-1}}
\newcommand{\ot}{\leftarrow}
\newcommand{\ninfty}{-\infty}
\newcommand{\pinfty}{+\infty}
\newcommand{\floor}[1]{\left\lfloor #1 \right\rfloor}
\newcommand{\ceil}[1]{\left\lceil #1 \right\rceil}
\newcommand{\viff}{\Big\Updownarrow} % vertical iff
\newcommand{\bmid}{\,\middle\vert\,} % big mid

% probability
\newcommand{\Prob}{\mathbf{P}}
\renewcommand{\vec}[1]{\boldsymbol{\mathbf{#1}}}
\DeclareMathOperator{\Bin}{Bin}
\DeclareMathOperator{\Geo}{Geo}
\DeclareMathOperator{\Poi}{Poi}
\DeclareMathOperator{\Exp}{Exp}
\DeclareMathOperator{\Var}{Var} % Variance
\DeclareMathOperator{\Cov}{Cov}

% special letters
\newcommand{\N}{\mathbb{N}}
\newcommand{\Z}{\mathbb{Z}}
\newcommand{\Q}{\mathbb{Q}}
\newcommand{\R}{\mathbb{R}}
\newcommand{\C}{\mathbb{C}}
\newcommand{\F}{\mathbb{F}}
\newcommand{\E}{\mathbf{E}}
\newcommand{\D}{\mathbb{D}}
\renewcommand{\H}{\mathbb{H}}
\newcommand{\ps}{\mathcal{P}}
\newcommand{\M}{\mathcal{M}}
\renewcommand{\L}{\mathcal{L}}
\renewcommand{\O}{\mathcal{O}}
\renewcommand{\U}{\mathcal{U}}
\newcommand{\Omicron}{O}
\newcommand{\powerset}{\mathcal{P}}

% text
\newcommand{\st}{\text{ s.t. }}
\newcommand{\tand}{\quad \text{and} \quad}
\newcommand{\tor}{\quad \text{or} \quad}
\newcommand{\stand}{\text{ and }}
\newcommand{\stor}{\text{ or }}
\renewcommand{\tt}[1]{\textnormal{\textbf{(#1).}}} %tt=theorem title GET RID OF

% title format
\title{\textbf{\ncourse}}

\ifx \ntype \notes
  \author{Notes taken by \nauthor\, \nemail \\ Based on lectures by \nlecturer}
  \date{\nterm\ \nyear}
\else
  \ifx \ntype \problems
    \author{These problems were solved by \nauthor \\ \nemail}
  \else
  \fi
\fi

\makeatother

\DeclareMathOperator{\gon}{gon}


\begin{document}
\maketitle

% Insert cool image here

\newpage
\tableofcontents
\newpage

\section{Introduction}

\subsection{Algebraic subsets of \texorpdfstring{$\C^n$}{C} and Hilbert's nullstellensatz}

Algebraic geometry deals with sets of solutions to systems of polynomial
equations.

\begin{example}
  The solution set for $xy - 1 = 0$ in $\C^2$.
\end{example}
\begin{example}
  The solution set for $x^n + y^n = 1$ in $\Q^2$ for $n \geq 3$.
  By Fermat's last theorem there are no solutions to this equation
  unless either $x$ or $y$ is zero.
\end{example}
\begin{example}
  The solution set for the following system of polynomials:
  \[
    \begin{cases}
      xt - yz = 0 \\
      xt - y^2 = 0 \\
      yt - z^2 = 0
    \end{cases}
  \]
  in $\C^4$.
\end{example}

\begin{example}
  Consider an $n \times n$ matrix $A$ as a point in $\C^{n^2}$.
  Then we have that
  \[
    \set{A \mid \tr(A) = 0} =
    \set{A \mid \sum_{i=1}^{n} a_{ii} = 0}
  \]
  is given as the solution set of a polynomial.
  Similarly
  \[
    \set{\det A = 0} = \set{\sum_{\tau \in S_n}\operatorname{sgn}(\tau)\,a_{1,\tau(1)}a_{2,\tau(2)} \ldots a_{n,\tau(n)} = 0},
  \]
  is given by the solution set for a polynomial in the entries
  of the matrix (by the formula with ). %
\end{example}

\begin{definition}[Set of zeros]
  For a subset $S \subseteq \C[x_1,\dots,x_n]$,
  we define its set of zeros by
  \[
    Z(S) :=
    \set{(a_1,\dots,a_n) \in \C^n \mid f(a_1,\dots,a_n) = 0 \quad \forall f \in S}.
  \]
\end{definition}

\begin{proposition}
  Here are some basic properties of sets of zeros:
  \begin{enumerate}
    \item[(1)] $Z(S) = Z\del{(S)}$ where $(S)$ is the ideal generated by $S$.
      Thus $Z(S) = Z(\set{f_1,\dots,f_m})$ for any set $\set{f_1,\dots,f_m}$
      such that $(f_1,\dots,f_m) = (S)$.
      Such a finite generating set exists by the Hilbert basis theorem,
      which tells us that $\C[x_1,\dots,x_n]$ is Noetherian.
    \item[(2)] If $I \subseteq J$ then $Z(I) \supseteq Z(J)$.
    \item[(3)] $Z(I + J) = Z(I) \cap Z(J)$,
    \item[(4)] $Z(I \cap J) = Z(IJ) = Z(IJ) = Z(I) \cap Z(J)$.
    \item[(5)] $Z(\sqrt{I}) = Z(I)$,
      we so can restrict to radical ideals.
  \end{enumerate}
\end{proposition}

\begin{remark}
  Recall that the radical of an ideal $I \normalin R$ is defined by
  \[
    \sqrt{I} :=
    \set{f \in \C[x_1,\dots,x_n] \mid f^k \in I \text{ for some } k \geq 1}.
  \]
\end{remark}

\begin{example}
  We have that $Z(\set{0}) = Z\del{(0)} = \C^n$,
  and $Z(\C[x_1,\dots,x_n]) = Z(\set{1}) = \emptyset$.
  Additionally,
  we have that $Z\del{(x_1-a_1,\dots,x_n-a_n)} = \set{(a_1,\dots,a_n)}$.
  Observe that $(x_1-a_1,\dots,x_n-a_n)$ is a maximal ideal.
\end{example}

\begin{definition}
  For a subset $V \subseteq \C^n$ we define the ideal:
  \[
    I(V) :=
    \set{f \in \C[x_1,\dots,x_n] \text{ such that $f(p) = 0$ for all $p \in V$}}.
  \]
  Equivalently,
  \[
    I(V) :=
    \bigcap_{p \in V} m_p,
  \]
  where $m_p = (x_1-a_1,\dots,x_n-a_n)$ is the maximal ideal corresponding
  to the point $p = (a_1,\dots,a_n)$.
\end{definition}

\begin{definition}[Zariski topology]
  The Zariski topology on $\C^n$ is defined by declaring the closed sets
  to be $Z(I)$.
\end{definition}

\begin{remark}
  Note that the Zariski topology is very weak,
  in the sense that it does not have a lot of open or closed sets.
  In particular, it is not Hausdorff.
\end{remark}

\begin{question}
  For $n = 1$ what are proper closed subsets of $\C$?
\end{question}
\begin{answer}
  If $0 \neq I \properideal \C[x]$,
  then since $\C[x]$ is a PID and algebraically closed $I = (f(x))$
  for some polynomial $f(x) = \prod (x - a_i)$.
  Thus the Zariski closed subsets of $\C$ are precisely the finite sets.
  This implies that in this case,
  the Zariski topology agrees with the cofinite topology.
\end{answer}

\begin{theorem}[Weak Hilbert's nullstellensatz]
  All maximal ideal of $\C[x_1,\dots,x_n]$ are of the form
  $(x_1 - a_1,x_2 - a_2,\dots,x_n - a_n)$ for some $(a_1,\dots,a_n) \in \C^n$.
\end{theorem}
\begin{corollary}
  \label{cor:zorns}
  Let $I \idealin \C[x_1,\dots,x_n]$.
  Then $Z(I) = \emptyset$ if and only if $I = (1) = \C[x_1,\dots,x_n]$.
\end{corollary}
\begin{proof}
  We have seen that if $I = (1)$ then $Z(I) = \emptyset$.
  If $I \idealin \C[x_1,\dots,x_n]$ %not equal
  then by applying Zorn's lemma $I \subseteq m_p$
  for some maximal ideal $m_p$, which is a point ideal for some point
  $p = (a_1,\dots,a_n)$ by the weak nullstellensatz theorem,
  thus $p \in Z(I)$.
  This implies that $Z(I) \neq \emptyset$,
  which completes the proof.
\end{proof}

\begin{remark}
  As a consequence of the weak nullstellensatz theorem, the set
  $Z(I)$ can be identified with the set of maximal ideals of $C[x_1,\dots,x_n]$
  which contain $I$, which in turn are in natural bijection to the
  set of maximal ideals of the quotient ring of functions
  $k[X] := C[x_1,\dots,x_n]/I$ from the correspondence theorem.
\end{remark}

\begin{theorem}[Regular strength Hilbert's nullstellensatz]
  \label{thm:hilbert-nullstellensatz}
  For any ideal $I$ in $\C[x_1,\dots,x_n]$ we have
  \[
    I(Z(I)) = \sqrt{I}.
  \]
  That is, $I(Z(I))$ is the closure of $I$ with respect to the property
  of being radical.
\end{theorem}
\begin{proof}
  If $f \in \sqrt{I}$,
  then $f^k \in I$ for some $k \geq 1$.
  Therefore $f^k(p) = (f(p))^k = 0$ for all $p \in Z(I)$.
  This implies that $f(p) = 0$ for all $p \in Z(I)$.
  which implies that $f \in I(Z(I))$.

  Conversely (we are using Rabinowitch's trick, which is a particular case
  of localization).
  Consider $0 \neq f \in I(Z(I))$.
  In $\C[x_1,\dots,x_n,t]$ we claim that the ideal
  \[
    \hat I = (ft - 1) + (I) \subseteq \C[x_1,\dots,x_n,t]
  \]
  has no zeros, where $(I)$ denotes the ideal in the ring $\C[x_1,\dots,x_n,t]$
  generated by $I$.
  Indeed, for any $(a_1,\dots,a_n,b) \in Z(\hat I)$
  there must exist $g \in I$ such that $g(a_1,\dots,a_n) = 0$,
  and in particular $f(a_1,\dots,a_n) = 0$,
  so we get that $(ft - 1)$ has to vanish at $(a_1,\dots,a_n,b)$,
  or in other words that $0 = f(a_1,\dots,a_n) b - 1 = -1$
  which is a contradiction to $f(a_1,\dots,a_n) = 0$.
  
  By \Cref{cor:zorns}, the ideal $\hat I$ must be the whole
  ring, and in particular it contains $1 \in \hat I$.
  This means that there exist $f_i \in I$ such that:
  \[
    1 =
    \del{f(x_1,\dots,x_n) t - 1} g_0(x_1,\dots,x_n,t) +
    \sum_{i} f_i(x_1,\dots,x_n) g_i(x_1,\dots,x_n,t).
  \]
  Consider the image of this equality in the field of rational
  functions in $n$ variables $\C(x_1,\dots,x_n)$ under the algebra map
  which sends $x_i$ to $x_i$ and $t$ to $f(x_1,\dots,x_n)^{-1}$.
  Then the equation becomes:
  \[
    1 =
    0 +
    \sum_{i} f_i(x_1,\dots,x_n) \frac{h_i(x_1,\dots,x_n)}{f_i(x_1,\dots,x_n)^{k_i}},
  \]
  for some polynomials $h_i$ and some $k_i \geq 0$.
  After clearing the denominators we conclude that $f^k$ lies in $I$,
  which proves the other side of the inclusion and concludes the proof.
\end{proof}

\begin{remark}
  \Cref{thm:hilbert-nullstellensatz} says that there is a one-to-one
  correspondence between closed algebraic subsets of $\C^n$ (i.e. the $Z(I)$)
  and radical ideals of $\C[x_1,\dots,x_n]$.
\end{remark}

\subsection{Irreducible components,  dimension theory.
  Projective spaces. Algebraic varieties}

\begin{definition}[Reducible zero locus]
  $Z(I)$ is reducible if it can be written as a union of two distinct
  proper closed subsets $Z(I) = Z(I_1) \cup (I_2)$.
  It is called irreducible otherwise.
\end{definition}

\begin{example}
  For example when $n = 2$ and $I = (x_1,x_2)$
  we can write $Z(I) = Z(x_1) \cup Z(x_2)$.
\end{example}

\begin{remark}
  Disconnected algebraic sets are always reducible,
  however some reducible algebraic sets are not disconnected.
  You may try to think of some examples by yourselves, but some nice
  examples are also written in the corresponding part in Lev's notes.
\end{remark}

\begin{theorem}
  Every closed subset $Z(I)$ can be uniquely written as a union of irreducible
  closed subsets of $Z(I)$ none of which is contained in another.
\end{theorem}
\begin{proof}
  First we will prove existence.
  Assume, by way of contradcition that the claim is false for some $Z(I)$.
  Then we can keep decomposing it into smaller unions of smaller and smaller
  subsets indefinitely.
  This indefinitely decreasing sequence of closed algebraic subsets
  correspond to an indefinitely increasing sequence of ideals in
  $C[x_1,\dots,x_n]$ in contradiction to it being Noetherian.
  This implies we can write $Z(I)$ as a finite union of irreduible subsets,
  and it remains to remove subsets that are contained in others.

  In the uniqueness direction, suppose we have two finite decompositions
  as described in the theorem:
  \[
    Z(I) =
    \bigcup_{\alpha} Z(I_{\alpha}) =
    \bigcup_{\beta} Z(I_{\beta}).
  \]
  For each $\alpha$ we have that
  \[
    Z(I_{\alpha}) =
    \bigcup_{\beta} \del{Z(I_{\beta}) \cap Z(I_{\alpha})}.
  \]
  By irreducibility of $Z(I_{\alpha})$ one of $Z(I_{\beta}) \cap Z(I_{\alpha})$
  must equal $Z(I_{\alpha})$.
  This implies that for every $\alpha$ there exists $\beta$ such that
  $Z(I_{\alpha}) \subseteq Z(I_{\beta})$.
  Of course, the same is true for each $\beta$,
  so we get that $Z(I_{\alpha}) \subseteq Z(I_{\beta}) \subseteq Z(I_{\alpha'})$,
  and thus $Z(I_{\alpha}) \subseteq Z(I_{\alpha'})$, which implies
  that $\alpha = \alpha'$ and thus $Z(I_{\alpha}) = Z(I_{\beta})$.
  Since this works for each $\alpha$ and $\beta$ we have proved uniqueness.
\end{proof}

\begin{remark}
  The decomposition of $Z(I)$ into irredubile components is a particular example
  of the commutative algebra statement that for any ideal $I$
  in a Noetherian ring, the radical $\sqrt{I}$ can be uniquely written as:
  \[
    \sqrt{I} = \bigcap_{i} p_i,
  \]
  where $p_i$ are a finite set of minimal primes that contain $I$.
\end{remark}

\begin{definition}[Algebraic subvariety]
  If $Z(I)$ is irreducible, $Z(I)$ is called an algebraic subvariety of $\C^n$.
\end{definition}

\begin{remark}
  We note that algberaic subvarieties $Z(I)$ are very similar (and indeed
  not very far from) smooth complex manifolds.
  We have a concept of dimension for complex manifolds,
  so it is natural to expect a notion of dimension on $Z(I)$ as well.
\end{remark}

\begin{definition}[Dimension of subvarieties]
  For an irreducible $Z(I)$ the ideal $I$ is prime (by exercise),
  so we can look at the quotient ring $A = \C[x_1,\dots,x_n] / I$
  which is an integral domain.
  Its fraction field $\Frac(A)$ is a finitely generated field extension of $\C$,
  and we define:
  \[
    \dim_{\C} Z(I) := \trdeg_{\C} \Frac(A),
  \]
  where $\trdeg_{\C} \Frac(A)$
  represents the transcendence degree of $\Frac(A)/\C$.
\end{definition}

\begin{definition}[Krull dimension]
  For a topological space $X$ the Krull dimension of $X$ is $\sup n$ where
  $Z_0 \subsetneq Z_1 \subsetneq \cdots \subsetneq Z_n$
  is an increasing chain of irreducible closed subsets of $X$.
  Similarly, the Krull dimension of a ring $R$
  is the supremum over the length of all increasing
  chains of prime ideals.
\end{definition}

\begin{remark}
  Note that $\Frac(\C[x_1,\dots,x_n] / I)$ is simply the fraction field
  of the coordinate ring $k[x_1,\dots,x_n]$, which are the functions over $X$
  (we have seen this in commutative algebra).
  Similarly, we intuitively look at $\Frac(\C[x_1,\dots,x_n] / I)$
  as the field of ``rational functions on $X = Z(I)$''.
\end{remark}

\begin{example}
  Clearly, the dimension of $\C^n$ is the transcenence degree of the field
  of rational functions $C(x_1,\dots,x_n)$ which is $n$.
\end{example}

\begin{remark}
  Any chain of irreducible closed subsets in $\C^n$
  can be completed to a chain with a member in each dimension.
  The corresponding property for chains of prime ideals in $\C[x_1,\dots,x_n]$
  is called \textit{catenary}.
  Rings with this property are called catenary rings.
\end{remark}

\begin{remark}
  Algebraic geometers have a very annoying (for beginners) of drawing complex
  varieties of dimension $n$ as objects of real dimension $n$,
  with the execption of Riemann surfaces (which are of complex dimension $1$),
  and are sometimes drawn as real surfaces.
\end{remark}

\begin{definition}[Complex projective space]
  The complex projective space denoted $\CP^n$ is defined as the set
  \[
    \del{\C^{n + 1} \setminus \set{\vec 0}} / \sim
  \]
  where $\vec a = (a_0,\dots,a_n) \sim \vec b = (b_0,\dots,b_n)$ if and only if
  there exists $\lambda \in \C^{\times}$ such that $\vec a = \lambda \vec b$.
\end{definition}
\begin{remark}
  The points in $\CP^n$ are written $(a_0 : a_1 : \cdots : a_n)$
  in order to distinguish them from points in $\C^{n+1}$.
\end{remark}

\begin{remark}
  We can think of $\CP^n$ as all the lines passing through the origin
  of $\C^{n+1}$.
\end{remark}

\begin{remark}
  For every $i \in \set{0,1,\dots,n}$ we define the subspace
  \[
    U_i = \set{(x_0 : \cdots : x_i : \cdots : x_n) \in \CP^n \mid x_i \neq 0}.
  \]
  By the structure of $\CP^n$, these points can be normalized so that:
  \[
    U_i =
    \set{\left(\frac{x_0}{x_i} : \cdots : 1 : \cdots : \frac{x_n}{x_i}\right)
      \in \CP^n} \cong \C^n.
  \]
  One can use the covering $\set{U_i}$ of $\CP^n$ to give it
  the structure of a complex manifold.
\end{remark}

\begin{definition}[Algebraic variety, rough definition]
  An algebraic variety is a topological space $X$ with a finite covering by
  $U_1,\dots,U_m$
  open subsets such that $U_i \cong Z(I_i) \subset \C^{j_i}$
\end{definition}

\begin{remark}
  We can show that two different closed subvarieties $Z(I) \subseteq \C^n$
  and $Z(J) \subseteq \C^m$ are isomorphic as algebraic varieties
  if and only if the corresponding finitely generated coordinate rings
  (generated algebras) $\C[x_1,\dots,x_n]/I$ and $\C[x_1,\dots,x_m]/J$
  are isomorphic.
  It is a natural statement, since points are in bijections with
  maximal ideals in these rings. We think of them as just the same variety in
  two different realizations. Such varieties are called \textit{affine}.
\end{remark}

\begin{remark}
  The main adventage of $\CP^n$ over $\C^{n}$ is that it is compact.
  Intuitively, it is the ``smallest'' way to compactify $\C^n$ to a compact
  complex manifold.
  In contrast, if we just wanted a compactification as a real manifold,
  we could do it with adding a single point.
\end{remark}

\begin{theorem}
  $\CP^n$ is a compact complex manifold.
\end{theorem}
\begin{proof}
  We can take the quotient in two steps.
  First by $\R_{> 0}$, then by $S^1 = \set{z \in \C^{\times} \mid |z| = 1}$.
  We immediately see that (as real manifolds):
  \[
    \del{\C^{n+1} \setminus \set{\vec 0}} / \R_{> 0} \cong
    \set{(x_0,\dots,x_n) \in \C^{n+1} \mid |x_0|^2 + \cdots + |x_n|^2 = 1}
    \cong S^{2n + 1}.
  \]
  Then $\CP^n$ is the image of the compact set $S^{2n + 1}$ under a continuous
  map and is therefore compact.
\end{proof}

\begin{definition}
  Consider $\C[x_0,\dots,x_n]$.
  Notice that this is a graded ring $\C[x_0,\dots,x_n] = \bigoplus_{d \geq 0} S_d$,
  where:
  \[
    S_d := \vspan\set{x_0^{a_0} \cdots x_n^{a_n} \mid \sum a_i = d},
  \]
  which are the homogeneous polynomials of degree $d$.
\end{definition}

\begin{definition}
  An ideal $I \idealin \C[x_0,\dots,x_n]$ is called homogeneous
  if it is generated by homogeneous polynomials (recall $\C[x_0,\dots,x_n]$
  is Noetherian and thus generated by finitely many polynomials).

  For a homogeneous ideal $J \idealin \C[x_0,\dots,x_n]$ we define:
  \[
    Z_+(J) := \set{(a_0 : \cdots : a_n) \in \CP^n \mid f(a_0,\dots,a_n) = 0
    \quad \forall \text{homogeneous } f \in J},
  \]
  which is well-defined since $f$ indeed vanishes on all of $(a_0 : \cdots : a_n)$
  since it is homogeneous.
  These are called projective algebraic subsets.

  Similarly, we define for $V \subset \CP^n$ the set:
  \[
    I_+(V) := \left(\text{homogeneous } f \in \C[x-0,\dots,x_n] \mid f(p) = 0 \quad \forall p \in V\right),
  \]
  which is a homogeneous ideal since it is generated by homogeneous ideals.
\end{definition}

\begin{remark}
  The Zariski topology on $\CP^n$ is defined by declaring the closed subsets
  to be the $Z_+(J)$.
\end{remark}

\begin{remark}
  We saw that $\CP^n = \bigcup_{i=0}^{n} U_i$ where $U_i \cong \C^n$.
\end{remark}

\begin{question}
  Is the topology induced from the Zariski topology on $\CP^n$ the same
  Zariski topology we defined on $\C^n$?
\end{question}
\begin{answer}
  Yes, because $U_i \cap Z_+(J) = Z(\widetilde T^i)$ where
  \[
    \widetilde J^i = \langle f(\frac{x_0}{x_i},\dots,\frac{x_n}{x_i} \in
    \C[\frac{x_0}{x_i},\dots,\frac{x_n}{x_i}] \mid \forall f \in J \text{ homogeneous})\rangle
  \]
  This process is called dehomogenization.
\end{answer}

\begin{question}
  What is the relation between $Z_+(J) \subset \CP^n$ and $Z(J) \subset \C^{n+1}$?
  The set $Z(J)$ is called the affine cone over $Z_+(J)$.
  % visualize this! lines in CP^2
\end{question}

\begin{remark}
  Relationship between closed subsets of $\CP^n$ and ideals is a little more
  complicated.
\end{remark}

\begin{remark}
  If $J$ is homogeneous then $J = \bigoplus_{d \geq 0} J_d$
  and $Z_+(J) = Z_+(J')$ if there exists $d_0$ such that $J_d = J'_d$
  for all $d \geq d_0$. (throwing out the first few terms).
\end{remark}

\begin{example}[Irrelevant ideal]
  The ideal $(x_0,x_1,\dots,x_n)$ is called the irrelevant ideal.
  This is the unique homogeneous maximal ideal.
\end{example}

\begin{example}
  Consider $\set{x_0 x_2 - x_1^2 = 0} \subset \CP^2$.
  If we dehomogenize with respect to $x_0$ we set it to $1$
  and get:
  \[
    \bigcap U_0 \colon \set{y_2 - y_1^2} \subset \C^2.
  \]
\end{example}

\begin{example}
  Smooth quadric degree 2 homogeneous polynomial
  $\set{x_0 x_3 - x_1 x_2 = 0} \subset \CP^3$.
  We will show that this is isomorphic to $\CP^1 \times \CP^1$.
\end{example}

\begin{example}
  $\set{y^2 - 4x^3 - g_2 x z^2 - g_3 z^3 = 0} \subset \CP^2$
  with coordinates $(x : y : z)$.
  This is known as an elliptic curve.
\end{example}

\section{Bezout's theorem}

\begin{remark}[Manifistation of compactness]
  Let $I$ be a homogeneous ideal.
  Then we can define:
  \[
    a_d := \dim_{\C}\left(\C[x_0,\dots,x_n] / I\right)_{\deg = d} < \infty.
  \]
\end{remark}

\begin{theorem}
  For a nontrivial (not irrelevant for example)
  homogeneous ideal $I \idealin \C[x_0,\dots,x_n]$,
  the coefficients $a_d$ satisfy:
  \begin{enumerate}
    \item[(1)] There exists a polynomials $P(z) \in \Q[z]$ such that $P(d) = a_d$
      for all sufficiently large $d > 0$.
      The polynomial $P$ is called the \textit{Hilbert polynomial} of $Z_+(I)$.
    \item[(2)] $\deg(P) = \dim Z_+(I)$.
    \item[(3)] $(\dim Z_+(I)) \mid (\text{leading coefficient of } P) \in \N$,
      which is called the degree of $Z_+(I)$.
    % spoiler, this is = to the number of points in Z_+(I) \cap H_1 \cap \cdots \cap  H_{\dim Z_+(I)}
    % where $H_i$ is a gen copy of CP^{n-1} \subset \CP^n.
    \item[(4)] The Hilbert function $f(t) := \sum_{d \geq 0} a_d t^d$ is a rational
      function of the form
      \[
        f(t) = \frac{g(t)}{(1 - t)^{\dim Z(I) + 1}},
      \]
      with a polynomial $g(t)$ such that $g(1) \neq 0$.
  \end{enumerate}
\end{theorem}

\begin{example}
  Take the homogeneous ideal $I = \set{0}$ (which corresponds to the whole space).
  We have that $a_d$ is the dimension of the spce of deg $d$ of homogeneous
  polynomials in $n + 1$ variables.
  This is equal to the number of ways of picking $n$ seperators among $n + d$
  dots.
  Thus, we get for all $d \geq 0$
  \[
    a_d := \binom{n+d}{n} = \frac{1}{n!} (d + n) \cdots (d + 1) =
    \frac{1}{n!} d^n + \text{etc.},
  \]
  times dimension factorial the leading coefficient is $1$,
  and this is indeed the dimension which we expect.

  The generating function:
  \[
    \sum_{d \geq 0} \frac{t^d}{n!} (d + n) \cdots (d + 1) =
    \sum_{d \geq -n} \frac{t^d}{n!} (d + n) \cdots (d + 1) =
    \frac{1}{n!} \del{\frac{\partial}{\partial t}}^n \del{\frac{1}{1 - t}} =
    \frac{1}{(1 - t)^{n+1}},
  \]
  so the claims of the previous theorem hold.
\end{example}

\begin{example}
  Consider a curve $C \subset \CP^2$ given by $I = (f)$,
  for a homogeneous $f \in \C[x_0,x_1,x_2]$ of degree $d$ without repeated
  factors.

  We have that:
  \[
    \dim\left( \C[x_0,x_1,x_2] / I\right)_{\deg = d} =
    \dim\left( \C[x_0,x_1,x_2]\right)_{\deg = d} -
    \dim\left( \C[x_0,x_1,x_2]\right)_{\deg = d - k}
  \]
  because of $xf \colon \C[x_0,x_1,x_2] \to \C[x_0,x_1,x_2]$ is injective.
  Thus,
  \[
    a_d =
    \begin{cases}
      \frac{(d + 2)(d + 1)}{2} - \frac{(d - k + 2)(d - k + 1)}{2}, & d \geq k; \\
      \frac{(d + 2)(d + 1)}{2}, & 0 \le d < k.
    \end{cases}
  \]
  Thus, for sufficiently large $d$ we get
  \[
    a_d = dk - \frac{k (k - 3)}{2},
  \]
  which is linear in $d$.
  We get that $\deg Z_+(I) = k$.
  And the generating function is:
  \[
    f(t) := \frac{1}{(1 - t)^3} - \frac{t^k}{(1 - t)^3} =
    \frac{\sum_{i=0}^{k-1} t^i}{(1 - t)^2},
  \]
  so again the last theorem holds.
\end{example}

% interlude on graded modules

\begin{definition}
  A $\bigoplus S_d = \C[x_0,\dots,x_n]$-module $M$ is called a graded module
  if $M = \bigoplus_{d \in \Z} M_d$,
  where $M_d$ is an abelian group such that $S_e \cdot M_d \subset M_{d + e}$.
\end{definition}

\begin{definition}[Shifted module]
  The shifted module $M(r)$ has the same module strucure but with shifted
  grading:
  \[
    M(r)_{\deg = k} := M_{\deg = k + r}.
  \]
  For any graded module $M$ we get a (degree $0$) graded map
  of graded modules $M(-1) \taking{x_n} M$ (multiplication by $x_n$?)
  with kernel $K$ and cokernel $C$ (also graded modules).
  This implies that:
  \[\begin{tikzcd}[ampersand replacement=\&,cramped]
	0 \& K \& {M(-1)} \& M \& C \& 0
	\arrow[from=1-1, to=1-2]
	\arrow[from=1-2, to=1-3]
	\arrow[from=1-3, to=1-4]
	\arrow[from=1-4, to=1-5]
	\arrow[from=1-5, to=1-6]
\end{tikzcd}\]
  is an exact sequence of graded modules.
  Thus,
  \[\begin{tikzcd}[ampersand replacement=\&,cramped]
	0 \& K_d \& {M(-1)_d} \& M_d \& C_d \& 0
	\arrow[from=1-1, to=1-2]
	\arrow[from=1-2, to=1-3]
	\arrow[from=1-3, to=1-4]
	\arrow[from=1-4, to=1-5]
	\arrow[from=1-5, to=1-6]
\end{tikzcd}\]
  is an exact sequence of vector spaces for all $d \geq 0$,
  and thus
  \[
    -\dim K_d + \dim M(-1)_d + \dim M_d + \dim C_d = 0.
  \]
  This happens if and only if
  \[
    -f_k(t) + f_{M(-1)}(t) - f_M(t) + f_C(t) = 0,
  \]
  then we observe that $f_{M(-1)}(t) = t f_M(t)$ we get:
  \[
    (1 - t) f_M(t) = f_C(t) - f_K(t).
  \]
  Upshot: $\C[x_0,\dots,x_n]$-module section on $K$ and $C$ factors
  through $\C[x_0,\dots,x_{n-1}]$-module section $\C[x_0,\dots,x_n] / (x_n)$
  because $x_n$ goes to $0$ on $K$ and $C$.
  We now apply induction to get $(4)$.
  The proof of $(1)$-$(3)$ is significantly more difficult and is this omitted.
\end{definition}

\begin{remark}
  We defined $\dim Z(I)$
  for irreducible $Z(I)$ % as $\trdeg \Frac(\C[x_1,\dots,x_n]/I)$.

  By analogy we get for irreducible $Z_+(I)$
  that:
  \[
    \dim Z_+(I) \cap U_0 = \trdeg \Frac(\C[x_1,\dots,x_n]/\widetilde I).
  \]
  Maybe on any other $U_i$, say $U_1$, we get a different definition? NO!
  Because % to be continued
\end{remark}

\begin{remark}
  Expressing $\Frac(A) := \Frac(\C[x_1,\dots,x_n]/\widetilde I)$
  intrinsically in terms of the variables $x_1,\dots,x_n$ we see that:
  \[
    \Rat(\Z_+(I)) = \set{\frac{g_1}{g_2}, \deg g_1 = \deg g_2, g_1,g_2 \in A}
    \subset \Frac(A).
  \]
\end{remark}

\begin{theorem}[Bezout's theorem]
  Let $C_1,C_2 \subset \CP^2$ be curves with no common components.
  Then the number of points in $C_1 \cap C_2$ counting multiplicities in
  $(\deg  C_1)(\deg  C_2)$.
\end{theorem}

% krulls principle ideal theorem | height ideals? - codimension 1 ideals
% every curve C \subset \CP^@ in Z_+(O) for I=(f)
% f homog is a polynomial in C[x_0,x_1,x_2].

\begin{remark}
  Suppose that $C_1$ nd $C_2$ intersect at $(1:0:0) \in \CP^2$
  and $(1:0:0) \in U_0 \cong \C^2 \subset \CP^2$.
  Let $C_1 \cap U_0 = Z(f)$, $C_2 \cap U_2$ for some $f,g \in \C[y_1,y_2]$
  for $y_i = \frac{x_i}{x_0}$.

  Let
  \[
    R = \set{\frac{F(y_1,y_2)}{G(y_1,y_2)} \mid F,G \in \C[y_1,y_2] \st G(0,0) \neq 0}.
  \]
  % this is the localization with respect to...
  Because $f$ and $g$ have no irreducible components in common
  (since $C_1$ and $C_2$ have no common components).
  This implies that $\dim_G R/(f,g) < \infty$.
  We also define:
  \[
    \mathrm{mult}_{(1:0:0)} (\set{f=0} \cap \set{g=0}) :=
    \dim_G R/(f,g) < \infty
  \]
\end{remark}

\begin{example}
  Let $C_1 = Z_+(x_1)$ and $C_2 = Z_+(x_2)$.
  In $U_0$ we get $f = y_1$, $g = y_2$,
  and thus $R / (y_1,y_2) = \C$, % verify
  and so $\dim R/(y_1,y_2) = 1$,
  which is exactly the number of intersection the lines $C_1$ and $C_2$
  should have according to Bezout's theorem.
\end{example}

\begin{example}
  Let $C_1 = Z_+(x_2)$ and $C_2 = Z_+(x_0 x_2 - x_1^2)$.
  In $U_0$ we get $f = y_2$, $g = y_2 - y_1^2$,
  and then:
  \begin{align*}
    R / (f,g) &=
    (R / (f)) / (g) =
    \set{\frac{F(y_1)}{G(y_1)} \mid G(0,0) \neq 0} / (y_2 - y_1^2) \\ &=
    \set{\frac{F(y_1)}{G(y_1)} \mid G(0,0) \neq 0} / (y_1^2) =
    \C \cdot 1 + \C \cdot y_1 \cong \C^2,
  \end{align*}
  so at $(1:0:0)$ the multiplicity of $C_1 \cap C_2$ is $2$.
\end{example}

\subsection{Proof of Bezout's theorem}
The idea of the proof is to consider $M = \C[x_0,x_1,x_2] / (f_1,f_2)$
as a graded $\C[x_0,x_1,x_2]$-module.

\begin{proposition}
  For every nonzero finitely generated graded module $M$ over a graded
  Noetherian ring $A$ there exists a finite filtration
  \[
    0 = M_0 \subsetneq M_1 \subsetneq \cdots \subsetneq M_k = M
  \]
  so that for each $i \in \set{1,\dots,k}$ the quotient module $M_i/M_{i-1}$
  is isomorphic to $(A/p_i)(d_i)$, where $p_i \idealin A$ is a homogeneous
  prime ideal in $A$ and $d_i \in \Z$ is a grading shift.
\end{proposition}
\begin{proof}
  When $M = \C[x_0,x_1,x_2] / (f_1,f_2)$,
  from the description of $M_i/M_{i-1}$ as $(A/p_i)(d_i)$
  we see $p_i = \Ann(m) := \set{a \in A \st am = 0}$.
  So any such $p \supseteq (f_1,f_2)$ if and only if
  $Z_+(p) \subset Z_+(f_1 f_2)$.
\end{proof}

\begin{remark}
  The Hilbert polynomial of any finitely generated
  graded $\C[x_0,x_1,x_2]$-module is $\sum_{i=1}^{h} HP\del{(A/p_i)(d_i)}$
  (the Hilbert polynomial of the summands).

  Note that $HP(\C[x_0,x_1,x_2]/(f_1,f_2))$ can be computed since:
%   \[\begin{tikzcd}[ampersand replacement=\&,cramped]
% 	0 \& {\left(\mathbb C[x_0,x_1,x_2]/(f_1)\right)(-\deg f_{2})} \& {\mathbb C[x_0,x_1,x_2]/(f_1)} \& {\mathbb C[x_0,x_1,x_2]/(f_1,f_2)} \& 0
% 	\arrow[from=1-1, to=1-2]
% 	\arrow["{f_2}", from=1-2, to=1-3]
% 	\arrow[from=1-3, to=1-4]
% 	\arrow[from=1-4, to=1-5]
% \end{tikzcd}\]
% https://q.uiver.app/#q=WzAsNSxbMCwwLCIwIl0sWzEsMCwiXFxsZWZ0KFxcbWF0aGJiIENbeF8wLHhfMSx4XzJdLyhmXzEpXFxyaWdodCkoLVxcZGVnIGZfezJ9KSJdLFsyLDAsIlxcbWF0aGJiIENbeF8wLHhfMSx4XzJdLyhmXzEpIl0sWzMsMCwiXFxtYXRoYmIgQ1t4XzAseF8xLHhfMl0vKGZfMSxmXzIpIl0sWzQsMCwiMCJdLFswLDFdLFsxLDIsImZfMiJdLFsyLDNdLFszLDRdXQ==
Note that the map labeled $f_2$ is injective since otherwise
$f_1$ and $f_2$ have a common factor.

We also label the terms in the sequence by $M_1(-\deg f_2)$, $M_1$
and $M$ respectively.

Note that as we have computed before $P_{M_1}(k) = (\deg f_1) k + c$
and $P_M(k) = (h \deg f_1 + c) - ((k - \deg f_2) \deg f_1 + 1) = \deg f_1 \deg f_2 = \deg C_1 \deg C_2$.

% this completes the first side

Next consider the irrelevant ideal $P = (x_0,x_1,x_2)$.
We get that $\C[x_0,x_1,x_2] / P \cong \C$ so for $d = 0,1,2,\dots$
we get $a_d = 1,0,0,\dots$, so $P_{\Q[x_0,x_1,x_2]/P}(h) = 0$.

Next consider $P = (x_1,x_2)$ (the other cases are symmetrical).
We have the isomorphism $\C[x_0,x_1,x_2]/P \cong \C[x_0]$,
and thus for $d = 0,1,2,\dots$ we get $a_d = 1,1,1,\dots$
and therefore $P_{\Q[x_0,x_1,x_2]/P}(k) = 1$.

% two sentences picture on phone
\end{remark}

% detour on localization

Let $S \subset A$ (where $A$ is a ring) be a mult subsets i.e.
$1 \in S$ and $f,g \in S$ implies $fg \in S$.

The localization of $A$ at $S$, $A_s$
is the set $\set{\frac{a}{s} \mid a \in A \stand s \in S} / \sim$
where
\[
  \frac{a_1}{s_1} = 
  \frac{a_2}{s_2}
\]
if and only if there exists $s_3 \in S$ such that $s_3(s_1 a_2 - s_2 a_1) = 0$
in $A$.

% beautiful geometric interpretation of localization

Consider the map $\varphi \colon A \to As$ given by $a \mapsto \frac{a}{1}$.
For an $A$-module
\[
  M_S := \set{\frac{m}{s} \mid m \in M, s \in S}/\sim
\]
where
\[
  \frac{m_1}{s_1} = 
  \frac{m_2}{s_2}
\]
if and only if there exists $s_3 \in S$ such that $s_3(s_1 m_2 - s_2 m_1) = 0$
in $M$.

Define $\varphi_{M,S} \colon M \to M_s$ defined by $m \mapsto \frac{m}{1}$.

\begin{remark}[Fact]
  Localizing at $S$ is an exact functor.
  i.e. if
% https://q.uiver.app/#q=WzAsNSxbMCwwLCIwIl0sWzEsMCwiTSJdLFsyLDAsIk0nIl0sWzMsMCwiTScnIl0sWzQsMCwiMCJdLFswLDFdLFsxLDJdLFsyLDNdLFszLDRdXQ==
\[\begin{tikzcd}[ampersand replacement=\&,cramped]
	0 \& M \& {M'} \& {M''} \& 0
	\arrow[from=1-1, to=1-2]
	\arrow[from=1-2, to=1-3]
	\arrow[from=1-3, to=1-4]
	\arrow[from=1-4, to=1-5]
\end{tikzcd}\]
is a short exact sequence (SES) of $A$-modules,
then so is the corresponding short exact sequence of $A_S$-modules!
\[\begin{tikzcd}[ampersand replacement=\&,cramped]
	0 \& M_S \& {M_S'} \& {M_S''} \& 0
	\arrow[from=1-1, to=1-2]
	\arrow[from=1-2, to=1-3]
	\arrow[from=1-3, to=1-4]
	\arrow[from=1-4, to=1-5]
\end{tikzcd}\]
\end{remark}

\begin{remark}
  Suppose our $P = (x_1,x_2)$.
  Consider:
  \[
    S = \set{\text{homogeneous polynomials of $\C[x_0,x_1,x_2]$ that don't vanish at $(1:0:0)$}}
  \]
  Notice that $x_0 \in S$ so we can localize at $x_0$ first
  (i.e. take $T = \set{x_0^n \mid n \geq 0}$)
  and only then we divide $\C[x_0,x_1,x_2]_{x_0} := \C[x_0,x_1,x_2]/$ % complete, divided by what?
  Then we get $\C[\frac{x_1}{x_0},\frac{x_2}{x_0},x_0^{\pm 1}] = \C[y_1,y_2][x_0^{\pm 1}]$.
  Localizing all the way we get that $\Z$-graded ring $R[x_0^{\pm 1}]$
  where $R = \set{\frac{F(y_1,y_2)}{G(y_1,y_2)} \mid G(0,0) \neq 0}$.

  Consider $M_S \cong \del{R/(\tilde f_1, \tilde f_2)}[x_0^{\pm 1}]$
  where $\tilde f_i = \frac{f_i}{x_0^{\dy f_i}}$ the dehomogenization of the $f_i$,
  the degree $0$ part is $R/(f,g)$ from the definition of
\end{remark}

% lecture 3

\begin{corollary}
  Let $C \subset \CP^2$ be an irreducible $\deg C = d$ curve.
  Then $C$ admits a degree $d$ morphism onto $\CP^1$.
\end{corollary}
\begin{proof}
  Suppose $P := (1 : 0 : 0) \notin C$.
  Then the stereographic projection $\varphi \colon \CP^2 - \set{P} \to \CP^1$
  from $P$ given by $(x_0 : x_1 : x_@) \mapsto (x_1 : x_2)$.
\end{proof}

\begin{definition}[Gonality]
  The gonality of a curve $C$ is defined to be:
  \[
    \gon(C) :=
    \min\set{\deg f \mid f \colon C \to \RP^1 \text{ is nonconstant}}.
  \]
\end{definition}

\section{Line bundles and vector bundles}

\begin{remark}
  If $f \colon X \to \C^N$ is a holomorphic function on a compact complex
  manifold,
  then $f$ is constant.
  In particular this is true for $f \colon X \to \C$.
\end{remark}

% bounded function on C is one that can be extended holomorphically to CP^1

\begin{definition}[Vector bundle]
  A vector bundle of rank $r$ over $X$ is a complex manifold $W$
  with a holomorphic map $\pi \colon W \to X$ such that
  for every $p \in X$ there exists an open set $p \in U$
  and a biholomorphism $\varphi_U \colon U \times \C^r \to \pi^{-1}(U)$
  so that the diagram
  \[\begin{tikzcd}[ampersand replacement=\&,cramped]
    {U \times \C^r} \&\& {\pi^{-1}(U)} \\
    \\
    \& U
    \arrow[from=1-1, to=1-3]
    \arrow[from=1-1, to=3-2]
    \arrow[from=1-3, to=3-2]
  \end{tikzcd}\]
  commutes and is compatible with the complex vector space structures on the
  fibers.
\end{definition}

\begin{remark}
  A vector bundle of rank $1$ is called a line bundle.
\end{remark}

\begin{example}[Real vector bundle]
  The M\"obius strip.
  The trivial bundle and the infinite mobius strip over $S^1$.
\end{example}

\begin{definition}[Trivial vector bundle]
  A vector bundle $W$ of rank $r$ is called a trivial vector bundle
  if $W \cong X \times \C^r$ as a vector bundle.
  That is, there exists a biholomorphic map $\psi \colon W \to X \times \C^r$
  such that
  \[\begin{tikzcd}[ampersand replacement=\&,cramped]
    W \&\& {X \times \C^r} \\
    \\
    \& U
    \arrow[from=1-1, to=1-3]
    \arrow["\pi"', from=1-1, to=3-2]
    \arrow["{pr_1}", from=1-3, to=3-2]
  \end{tikzcd}\]
\end{definition}

\begin{remark}
  Everything you can do to vector spaces in linear algebra you can do to vector
  bundles.
\end{remark}

\begin{definition}[Picard group]

\end{definition}


% d-duple embedding / d-th verendese embedding




\end{document}
